% =============================================================================
%  §1 Weil--Deligne 表示（底本 p.2--9）
% =============================================================================
\section{Weil--Deligne 表示}

本节中固定一个非 Archimedes 局部域 $K$，其剩余特征为 $p$。记 $k$ 为其剩余域，
$q=\#k=p^{f}$。固定一个素元 $\varpi_K\in\mathcal{O}_K$。最后，设
$G_K=\Gal(K^{s}/K)$ 为固定的可分闭包 $K^{s}$ 在 $K$ 上的 Galois 群。

\subsection{Weil--Deligne 表示}

\subsubsection{Weil 群 $W_K$}

按定义我们有交换图
\[
\begin{array}{ccccccccc}
0 & \longrightarrow & I_K & \longrightarrow & W_K & \longrightarrow & \Frob^{\mathbb{Z}} & \longrightarrow & 0\\[2pt]
 & & \Big\downarrow\simeq & & \Big\downarrow & & \Big\downarrow\simeq & & \\[2pt]
0 & \longrightarrow & I_K & \longrightarrow & G_K & \longrightarrow & G_k & \longrightarrow & 0 .
\end{array}
\]
这里 $W_K$ 的拓扑取成使 $I_K\subset W_K$ 为开子群、并带上它通常的拓扑，而
$\Frob^{\mathbb{Z}}\simeq\mathbb{Z}$ 取离散拓扑。特别地，这\emph{不是}从 $G_K$ 诱导的子空间拓扑。

\begin{remark}
一般地，若 $K$ 是局部域或整体域，与 $K$ 相伴的 $W_K$ 的定义涉及一个三元组
$(W_K,\iota,\{r_E\})$：其中 $W_K$ 是拓扑群，$\iota\colon W_K\to G_K$ 是像稠密的连续同态，
而对每个有限扩张 $E/F$ 有 $r_E\colon C_E\xrightarrow{\ \sim\ }W^{\mathrm{ab}}_E$，
这里 $C_E$ 在局部情形等于 $E^{\times}$、在整体情形等于 $E^{\times}\backslash\mathbb{A}^{\times}_E$；
这些数据满足若干自然条件。可以证明这样的 $W_K$ 存在且唯一。在非 Archimedes 局部情形，
它与上面的定义一致：$\iota$ 是自然的包含，$r_E$ 是局部类域论给出的同构。
\end{remark}

\subsubsection{驯特征标}

固定 $K^{s}$ 中一族相容的本原单位根 $\{\zeta_n\}_{p\nmid n}$。于是有 $K$ 上的
最大非分歧扩张 $K^{\mathrm{nr}}=\bigcup_{p\nmid n}K(\zeta_n)$ 与最大驯分歧扩张
$K^{t}=\bigcup_{p\nmid n}K^{\mathrm{nr}}(\varpi_K^{1/n})$。野惯性群
$P_K=G(K^{s}/K^{t})$ 是 $I_K$ 中唯一的 pro-$p$ Sylow 子群，而\emph{驯特征标}
（依赖于 $\{\zeta_n\}$ 的选取）是如下同构：
\[
t\colon I_K/P_K\xrightarrow{\ \sim\ }\prod_{l\neq p}\mathbb{Z}_l,
\qquad
\sigma\longmapsto t(\sigma)\ \mathrm{mod}\ n
\quad\text{其中}\quad
\frac{\sigma(\varpi_K^{1/n})}{\varpi_K^{1/n}}=\zeta_n^{t(\sigma)} .
\]
此外，对 $l\neq p$，令 $K^{t,l}=K^{\mathrm{nr}}(\varpi_K^{1/l^{\infty}})$ 为最大驯分歧扩张的
$l$-部分，则有（作为 Abel 群的同构）
\[
t_l\colon I_{K,l}/P_K=G(K^{t,l}/K^{\mathrm{nr}})\xrightarrow{\ \sim\ }\mathbb{Z}_l ,
\]
它与驯特征标到 $\mathbb{Z}_l$ 的投影一致。

$G_K$（或 $W_K$）通过共轭作用在 $I_K/P_K$ 上（因为 $P_K$ 是唯一的 pro-$p$ Sylow 子群）。
在每个 $\mathbb{Z}_l$ 分量上，该作用由 $l$-进分圆特征标 $\chi=\chi_l$ 给出：
\[
I_{K,l}/P_K\xrightarrow{\ \sim\ }\mathbb{Z}_l(1),
\]
其中的 $G_K$-作用是共轭。换言之，对任意 $\sigma\in G_K$、$\tau\in I_{K,l}$，有
\[
t_l(\sigma\tau\sigma^{-1})=\chi(\sigma)\,t_l(\tau).
\]

\begin{yizhu}[$\mathbb{Z}_l(1)$ 的读法]
底本把 $I_{K,l}/P_K\cong\mathbb{Z}_l$ 与 $\mathbb{Z}_l(1)$ 等同，靠的正是「共轭作用 =
乘 $\chi_l$」。因此这里的 $\mathbb{Z}_l(1)$ 不是抽象群，而是带 $G_K$-作用的模：
$G_K$ 通过 $\chi_l$ 作用在 $\mathbb{Z}_l$ 上。
\end{yizhu}

\subsubsection{Weil--Deligne 表示}

设 $L$ 是特征 $0$ 的域，$V$ 是 $L$ 上的有限维向量空间。

\begin{definition}
$W_K$ 在 $V$ 上的一个\emph{Weil--Deligne 表示}是一对 $(\rho,N)$，其中
\begin{enumerate}[label=(\arabic*),leftmargin=2.2em]
\item $\rho$ 是 $W_K$ 在 $V$ 上的连续表示（$L$ 取离散拓扑），等价地，$I_K$ 的像有限；
\item $N\in\End(V)$ 是自同态，使得对任意 $\sigma\in W_K$ 有
\[
\sigma N\sigma^{-1}=q^{-v_K(\sigma)}N ,
\]
其中 $q=\#k$，而 $v_K$ 由 $\sigma\mapsto\Frob_K^{v_K(\sigma)}\in W_K/I_K
\xrightarrow{\ \sim\ }\Frob^{\mathbb{Z}}_K$ 确定。
\end{enumerate}
\end{definition}

\begin{definition}
记 $\mathrm{wd}=\mathrm{wd}_L(W_K)$ 为 Weil--Deligne 表示构成的范畴，其态射是
与自同态 $N$ 相容的表示态射。
\end{definition}

\subsubsection{基本性质}

设 $(\rho,N)\colon W_K\to\GL(V)$ 是如上给出的 Weil--Deligne 表示。特别地，对任一
Frobenius 的提升 $\varphi\in W_K$，有
\[
\varphi N\varphi^{-1}=q^{-1}N .
\]
我们记录如下性质：

\begin{enumerate}[label=(\arabic*),leftmargin=2.2em]
\item $N$ 必然是幂零的：特征多项式满足
\[
f(X)=\det(I_n\cdot X-N)=\det(I_n\cdot X-q^{-1}N)=q^{-n}f(qX).
\]
于是 $f(X)=X^{n}$，这迫使 $N$ 幂零。

\item $N$ 显然不是 Galois 等变的；事实上 $N$ 诱导出一个 Galois 等变的映射
\[
N(-1)\colon V\longrightarrow V(-1),
\]
其中 $V(-1)$（作为一个形式符号）表示 $V$ 用非分歧特征标
$\epsilon\colon G_K\to L^{\times}$、$\varphi\mapsto q$ 所作的扭转，$\varphi$ 是 Frobenius 的一个提升。
若 $L/\mathbb{Q}_l$ 且 $l\neq p$，则 $(-1)$ 与 Tate 扭转的通常记法一致。

\item 设 $r\colon W_K\to\GL(V)$ 是 $W_K$ 的连续表示（$V$ 上取离散拓扑），则定义
$W_K$ 在 $V^{m}$ 上的 Weil--Deligne 表示 $\mathrm{Sp}_m(r)$ 为
\[
\mathrm{Sp}_m(r):=r\oplus r(1)\oplus\cdots\oplus r(m-1),
\]
其中 $r(i)$ 是 $r$ 用 $\chi^{-1}$ 所作的扭转；再定义 $N\in\End(V^{m})$ 使其对
$i\le m-2$ 把 $r(i)$ 同构地送到 $r(i+1)$，而在 $r(m-1)$ 上取 $N=0$。等价地，
\[
N(-1)\colon \mathrm{Sp}_m(r)=r\oplus\cdots\oplus r(m-1)\longrightarrow
\mathrm{Sp}_m(r)(-1)=r(-1)\oplus\cdots\oplus r(m-2)
\]
在 $r\oplus\cdots\oplus r(m-2)$ 上是恒等、在 $r(m-1)$ 上是 $0$。

\item 对使 $N^{m}=0$ 的最小 $m$ 作归纳，可以证明：每个使 $\rho$ 半单的
Weil--Deligne 表示 $(\rho,N)$ 都同构于若干形如 $\mathrm{Sp}_m(r)$ 的 Weil--Deligne 表示的直和，
其中 $r$ 是 $W_K$ 的不可约表示、且 $r(I_K)$ 有限。
\end{enumerate}

\begin{yizhu}[关于 $r(i)$ 的扭转方向]
底本写作「$r(i)$ 是 $r$ 用 $\chi^{-1}$ 所作的扭转」，而第 (2) 条又把 $V(-1)$ 定义为
用 $\epsilon(\varphi)=q$ 扭转。两处方向相反，读者容易在此处错一次符号。
译者按：本库统一采用「$V(1)=V\otimes\chi$，$V(-1)=V\otimes\chi^{-1}$」的 Tate 约定，
于是 $r(i)=r\otimes\chi^{-i}$ 与 $\mathrm{Sp}_m(r)$ 中 $N(-1)$ 把 $r(i)$ 送到 $r(i+1)$ 相容；
底本第 (2) 条的 $V(-1)$ 应读作「$N$ 把权 $-1$ 的 Tate 扭转送回权 $0$」这一形式陈述，
其符号以第 (3) 条的显式描述为准。
\end{yizhu}

\subsubsection{关于半单表示的注记}

本段讨论中设 $V$ 是任意域 $k$ 上的向量空间，$H<G$ 是子群，$\rho\colon G\to\GL(V)$ 是
$G$ 的表示。我们想知道 $\rho$ 的半单性与 $\rho|_H$ 的半单性之间的关系。

\begin{lemma}
\leavevmode
\begin{enumerate}[label=(\arabic*),leftmargin=2.2em]
\item 若 $V$ 有限维、$H\triangleleft G$ 是正规子群，则 $\rho$ 半单蕴含 $\rho|_H$ 半单。
\item 若 $H<G$ 的指数为 $r$，且 $r$ 在系数域 $k$ 中可逆，则 $\rho$ 半单当且仅当 $\rho|_H$ 半单。
\end{enumerate}
\end{lemma}

\begin{proof}
(1) 不妨设 $\rho$ 不可约。取 $W=\bigoplus_i W_i$，其中 $\{W_i\}$ 是 $V$ 的互不相同的
$H$-子表示，则由 $H$ 的正规性 $W$ 是 $V$ 的 $G$-子表示。由 $V$ 有限维知 $W\neq 0$，
于是 $W=V$，即 $\rho|_H=\bigoplus W_i$ 半单。

(2) 设 $G=\bigsqcup_{i=1,\dots,r}Hg_i$。

(a) 先设 $\rho$ 不可约。任一 $v\neq 0$ 作为单 $k[G]$-模生成 $V$，故 $\{g_iv\}$ 作为
$k[H]$-模生成 $V$。现取 $W\subset V$ 为极大的 $k[H]$-子模（因 $V$ 有限生成而存在）。
则每个 $V/g_iW$ 是单 $k[H]$-模，且有自然的 $k[H]$-线性映射 $V\to\bigoplus_i V/g_iW$，
其核是 $V$ 的真 $k[G]$-子模，故为 $0$。于是 $V\hookrightarrow\bigoplus_i V/g_iW$ 半单。

(b) 「若」的方向是 Maschke 定理。设 $\rho|_H$ 半单，$W\subset V$ 是 $G$-子表示，
从而也是半单的 $H$-子表示。取 $H$-线性投影 $\pi\colon V\to W$，对 $\pi$ 取平均得
$G$-线性投影
\[
\pi^{G}\colon V\longrightarrow W,\qquad
v\longmapsto\frac{1}{r}\sum_{i=1,\dots,r}g_i^{-1}\pi(g_i\cdot v),
\]
故得证。（要点在于：即使 $g_i$ 与 $h$ 不可交换，$g_i^{-1}\circ\pi\circ g_i$ 仍与 $h$ 交换，
因此 $\pi^{G}$ 是 $H$-线性的，从而是 $G$-线性的。）
\end{proof}

\subsubsection{$\varphi$-半单表示}

\begin{lemma}
设 $\rho$ 是 $W_K$ 在 $V$ 上的表示，且 $\rho(I_K)$ 有限。为简单起见设 $L$ 代数闭。则下列等价：
\begin{enumerate}[label=(\arabic*),leftmargin=2.2em]
\item 对一切 $\sigma\in W_K$，$\rho(\sigma)$ 半单；
\item 对 Frobenius 的（任一）提升 $\varphi$，$\rho(\varphi)$ 半单；
\item $\rho$ 是半单表示。
\end{enumerate}
\end{lemma}

\begin{proof}[略证]
注意 $\operatorname{char}L=0$。由于 $\operatorname{im}\rho(I_K)$ 有限，我们断言：对任意
$\sigma\in W_K$，存在 $m\in\mathbb{Z}_{>0}$ 使 $\rho(\sigma^{m})$ 落在 $\rho(W_K)$ 的中心里。
为证此断言，取 $\tau_1,\dots,\tau_r\in I_K$ 生成 $\rho(I_K)$，并令 $\tau_0=\varphi$ 是
Frobenius 的提升。对每个 $\tau_i$，考虑 $\rho(\tau_i^{n}\sigma\tau_i^{-n})$（$n\in\mathbb{Z}_{>0}$），
由 $\rho(I_K)$ 的有限性知存在 $m_i$ 使 $\rho(\tau_i^{m_i})$ 与 $\rho(\sigma)$ 交换；
取 $m=m_0m_1\cdots m_r$ 即可。

于是对任意 $\sigma\in W_K$，存在 $n\ge 0$、$n'>0$ 使 $\rho(\sigma^{n'})=\rho(\varphi^{n})$。
由于 $L$ 特征为 $0$，线性算子 $T$ 的半单性等价于 $T$ 的幂的半单性（注意这在正特征下
一般不成立），故 (1) 与 (2) 等价。又 $\rho(\varphi)$ 生成 $\rho(W_K)$ 中的有限指数子群，
而 $L$ 特征为 $0$，由引理 1.1 第 (2) 部分，$\rho(W_K)$ 在 $V$ 上半单作用当且仅当
$\langle\rho(\varphi)\rangle$ 半单作用，当且仅当 $\rho(\varphi)$ 半单。
\end{proof}

\begin{remark}
注意上一条引理中我们假定了 $L$ 代数闭。这是因为按我们的约定，线性算子称为半单，
指的是它「潜在可对角化」，即在某次基变换后可对角化。对一般的 $L$，只需把第 (3) 条
中的说法换成「潜在半单」。
\end{remark}

\begin{definition}
Weil--Deligne 表示 $(\rho,N)\in\mathrm{wd}_K$ 称为 \emph{$\varphi$-半单的}
（Frobenius 半单的），若 $\rho$ 半单；等价地，若 $\rho(\varphi)$ 半单。
\end{definition}

\begin{remark}
\leavevmode
\begin{enumerate}[label=(\arabic*),leftmargin=2.2em]
\item 设 $(\rho,N)\in\mathrm{wd}$。若 $\rho$ 不可约则 $N$ 必为 $0$，因为 $\ker N$
在 $W_K$ 下稳定。
\item $\varphi$-半单表示一般\emph{不是}范畴 $\mathrm{wd}$ 中的半单对象。
$\mathrm{wd}$ 中的半单对象是单对象的直和，即那些满足 $N=0$ 的 $\varphi$-半单表示。
\end{enumerate}
\end{remark}

\begin{remark}[Frobenius 半单化]
设 $(\rho,N)$ 是 $W_K$ 的 Weil--Deligne 表示，$\rho^{\mathrm{ss}}$ 是 $\rho$ 的半单化，
则 $(\rho^{\mathrm{ss}},N)$ 仍是 Weil--Deligne 表示，称为 $(\rho,N)$ 的
\emph{$\varphi$-半单化}。
\end{remark}

\subsubsection{张量}

设 $(\rho_1,N_1)$ 与 $(\rho_2,N_2)$ 是两个 Weil--Deligne 表示，则它们的张量积是
\[
(\rho_1\otimes\rho_2,\ N_1\otimes\mathrm{id}+\mathrm{id}\otimes N_2).
\]
类似地，$\Hom(\rho_1,\rho_2)=(r,N)$ 由下式给出：
\[
\bigl(r(\tau)\varphi\bigr)v_1=\rho_2(\tau)\cdot\varphi\bigl(\rho_1(\tau)^{-1}v_1\bigr),
\qquad
(N\varphi)v_1=N_2\cdot\varphi(v_1)-\varphi(N_1\cdot v_1).
\]

\subsubsection{Weil--Deligne 群概形}

把 $W_K$ 看作 $\mathbb{Q}$ 上的（常）群概形，它取得自极限
$\varprojlim_{I'\triangleleft I_K}W_K/I'$，其中 $I'$ 遍历 $I_K$ 的所有开正规子群。
这是 $\mathbb{Q}$ 上的无限型群概形。令
\[
\widehat{W}_K:=\mathbb{G}_a\rtimes W_K ,
\]
其中 $W_K$ 在 $\mathbb{G}_a$ 上的作用由 $wxw^{-1}=|w|_K x$ 给出。

\begin{center}
\emph{底本此处未写}（作者原注：to be completed in second draft）。
\end{center}

\subsubsection{Jacobson--Morozov}

\begin{center}
\emph{底本此处未写}（作者原注：to be completed in second draft）。
\end{center}

\subsubsection{Tate 与 Deligne--Langlands 的 $\epsilon$-因子}

我们需要为 Weil--Deligne 表示 $(\rho,N)$ 定义 $L$-导子与 $\epsilon$-导子，
在此过程中还要定义导子 $c((\rho,N))$。先对 $W_K$ 的表示做这件事。回忆：对我们而言，
$W_K$ 的表示指的是关于目标上离散拓扑连续的表示。

若固定一个同构 $L=\mathbb{Q}_l\xrightarrow{\ \sim\ }\mathbb{C}$，这些表示对应于
$W_K$ 的 Artin 表示。固定 $\varpi=\varpi_K$。

先回忆（Tate 论文中的结论）：对拟特征标 $\chi\colon K^{\times}\to\mathbb{C}$，Tate 定义局部
$L$-因子
\[
L(\chi)=
\begin{cases}
(1-\chi(\varpi))^{-1}, & \chi\ \text{非分歧},\\
1, & \chi\ \text{分歧}.
\end{cases}
\]
对 $\omega_s=\|\cdot\|^{s}$，$s\in\mathbb{C}$，$L(\chi\omega_s)$ 关于 $s$ 是亚纯的。
设 $\psi$ 是 $K$ 上非平凡加性特征标，$dx$ 是 $K$ 上的 Haar 测度。记 $d^{\times}x$ 为
$K^{\times}$ 上由 $\|x\|^{-1}dx$ 给出的 Haar 测度。最后，令
$\chi^{\vee}:=\chi^{-1}\omega_1=\chi^{-1}\|\cdot\|$ 为 $\chi$ 的扭对偶。
Tate 考虑 Fourier 变换 $\hat f(y)=\int f(x)\psi(xy)\,dx$，并证明
\[
\frac{\int\hat f^{\chi^{\vee}}(x)\,d^{\times}x}{L(\chi^{\vee})}
=\epsilon(\chi,\psi,dx)\cdot
\frac{\int f^{\chi}(x)\,d^{\times}x}{L(\chi)}
\]
其中 $\epsilon(\chi,\psi,dx)$ 是 $\chi$ 的非零全纯函数，称为 $\epsilon$-因子，
它依赖于 $\psi$ 与 $dx$。Tate 还证明：若令 $n(\psi)$ 为使 $\psi(\varpi^{-n}\mathcal{O}_K)=1$
的最大整数，$c(\chi)$ 为 $\chi$ 的导子，并取 $c\in K^{\times}$ 使 $v(c)=n(\psi)+c(\chi)$，则
\[
\epsilon(\chi,\psi,dx)=
\begin{cases}
\|c\|^{-1}\chi(c)\cdot\vol_{dx}(\mathcal{O}_K), & \chi\ \text{非分歧},\\[4pt]
\displaystyle\int_{c^{-1}\mathcal{O}_K^{\times}}\chi^{-1}(x)\psi(x)\,dx, & \chi\ \text{分歧}.
\end{cases}
\]
也就是说，他对特征标的 $\epsilon$-因子给出了显式公式。

对高维表示，（据 [10]）只有一条属于 Langlands 的存在性定理（Deligne 给了简短证明）。

\begin{theorem}
存在唯一的函数 $\epsilon$，对每个局部域 $K$、$W_K$ 的表示 $\rho$、非平凡加性特征标 $\psi$
与 $K$ 的 Haar 测度 $dx$，给出数 $\epsilon(\rho,\psi,dx)\in\mathbb{C}^{\times}$，使得
\begin{enumerate}[label=(\arabic*),leftmargin=2.2em]
\item 对特征标 $\chi$，$\epsilon(\chi,\psi,dx)$ 与上面的定义一致；
\item $\epsilon(\cdot,\psi,dx)$ 在 $W_K$ 表示的短正合列上是乘性的，特别地下降为
      Grothendieck 群上的映射 $\epsilon(\cdot,\psi,dx)\colon R(W_K):=K_0(\Rep W_K)\to\mathbb{C}^{\times}$；
\item $\epsilon$ 在下述意义下于 $0$ 次部分是归纳的：对每个域扩张 $E/L/K$ 与加性
      Haar 测度 $\nu_E$、$\mu_L$，有
\[
\epsilon(\Ind_{E/L}[r],\psi\circ\Tr_{L/K},\mu_L)
=\epsilon([r],\psi\circ\Tr_{E/K},\mu_E)
\]
对一切 $0$ 次元 $\ [r]\in R(W_K)$ 成立。
\end{enumerate}
\end{theorem}

\begin{remark}
设 $c(\rho)$ 是 $W_K$ 的表示 $(\rho,V)$ 的 Artin 导子，按定义等于
\[
c(\rho):=c(\chi_{\rho})=\sum_{i\ge 0}\frac{1}{[G_0:G_i]}\bigl(\dim V-\dim V^{G_i}\bigr),
\]
其中 $\chi_{\rho}$ 是该表示的特征标，$G_i$ 是 $I_K$ 的第 $i$ 个分歧子群。则可以证明
\[
\epsilon(\rho\cdot\omega_s,\psi,dx)=(\#k)^{-\bigl(c(\rho)+n(\psi)\dim V\bigr)\cdot s}\cdot
\epsilon(\rho,\psi,dx).
\]
\end{remark}

\subsubsection{导子、$L$-因子与 $\epsilon$-因子}

现在可以定义 Weil--Deligne 表示对应的不变量了。

\begin{definition}
设 $\rho=(\rho',N)$ 是 $W_K$ 的 Weil--Deligne 表示。设 $\psi$ 是 $K$ 上的加性特征标、
$dx$ 是 $K$ 上的加性 Haar 测度，$q=\#k$ 是剩余域的大小。
\begin{enumerate}[label=(\arabic*),leftmargin=2.2em]
\item 定义导子 $c(\rho)=c(\rho',N)$ 为
\[
c(\rho):=c(\rho')+\dim V^{I_K}-\dim(\ker N)^{I_K}.
\]
\item 设 $\Phi$ 是（任一）几何 Frobenius。定义 $L$-因子为
\[
L(\rho,s):=\det\Bigl(1-q^{-s}\Phi\ \big|\ (\ker N)^{I_K}\Bigr)^{-1}.
\]
\item 定义 $\epsilon$-因子为（下面省略了 $dx$ 的记号）
\[
\epsilon(\rho,s,\psi):=q^{-\bigl(c(\rho)+n(\psi)\dim V\bigr)\cdot s}\cdot
\det\Bigl(-\Phi\ \big|\ V^{I_K}/(\ker N)^{I_K}\Bigr)\cdot\epsilon(\rho',\psi,dx).
\]
\end{enumerate}
\end{definition}

特别地，从定义中挑出如下平凡的推论。

\begin{corollary}
\[
\epsilon(\rho,s,\psi)=q^{-\bigl(c(\rho)+n(\psi)\dim V\bigr)\cdot s}\cdot\epsilon(\rho,0,\psi).
\]
\end{corollary}

\begin{remark}
$\epsilon(\rho,s,\psi)$ 的定义与 [10] 略有不同：按我们的术语，Tate 把
$\epsilon$-因子定义为 $q^{n(\psi)\dim V\cdot s}\cdot\epsilon(\rho,s,\psi)$，
这似乎与上一页的注记不相容。我们的约定与 [11] 一致。
\end{remark}

\subsection{$l$-进表示}

\subsubsection{$l$-进 Galois 表示}

现在设 $L/\mathbb{Q}_l$ 是有限扩张，$l\neq p$。设 $V$ 是 $L$ 上的有限维向量空间。
固定一族相容的单位根 $\{\zeta_n\}_{p\nmid n}$ 与上面定义的 $l$-进驯特征标
$t_l\colon I_K/P_K\to\mathbb{Z}_l$。回忆 Grothendieck 的准单幂幺定理：

\begin{lemma}
设 $\rho\colon G_K\to\GL(V)$ 是 $G_K$ 的 $l$-进表示（即 $\rho$ 关于 $V$ 上的 $l$-进拓扑连续），则
\begin{enumerate}[label=(\arabic*),leftmargin=2.2em]
\item 存在开子群 $I'_K\subset I_K$（对应某个有限扩张 $K'/K$）与唯一确定的幂零算子
      $N\in\End(V)$，使得对一切 $\sigma\in I'_K$ 有 $\rho(\sigma)=\exp\bigl(t_l(\sigma)N\bigr)$；
\item 此外，对任意 $\tau=\sigma\varphi^{n}\in W_K$（其中 $\sigma\in I_K$，$\varphi$ 是 Frobenius 的
      固定提升），有 $\rho(\tau)N\rho(\tau^{-1})=q^{-n}N$。
\end{enumerate}
\end{lemma}

\begin{proof}[略证]
由连续性，$G_K$ 稳定某个 $\mathcal{O}_L$-格 $\Lambda\subset V$，模 $l$ 约化得
$\rho\colon G_K\to\Aut(\Lambda/l)$。其核 $\ker\rho=G_{K''}\subset G_K$ 是有限指数子群，
即 $K''/K$ 有限。由于 $\GL(\Lambda)$ 模 $l$ 约化的核是 pro-$l$ 的，知
$\rho|_{I_{K''}}$ 通过 $t_l\colon I_{K''}\to\mathbb{Z}_l(1)$ 分解。取 $\sigma\in I_{K''}$ 使
$t_l(\sigma)$ 拓扑生成 $\mathbb{Z}_l(1)$（作为 Abel 群）。

注意 $G_{K''}$ 在 $t_l(I_{K''})$ 上的共轭作用就是分圆特征标，即
$t_l(\varphi\sigma\varphi^{-1})=\chi(\varphi)t_l(\sigma)$，这迫使 $\rho(\sigma)$ 的特征值都是单位根。
于是存在有限扩张 $K'/K$ 使 $\rho(I'_K)$ 幂幺地作用。必要时把 $\sigma$ 换成 $I_{K'}$ 中新的
拓扑生成元（仍记作 $\sigma$），再令 $N=\log\rho(\sigma)$，它由幂幺性定义良好。
\end{proof}

\begin{corollary}
设 $\rho\colon G_K\to\GL(V)$ 是 $G_K$ 的 $l$-进表示。若 $\rho$ 半单，则 $\rho(I_K)$ 有限。
\end{corollary}

\begin{proof}
由引理 1.1 第 (1) 部分，$\rho|_{I_K}$ 半单。再把 $\rho$ 限制到引理 1.5 证明中的
$I_{K'}\triangleleft I_K$，则 $\rho|_{I_{K'}}$ 半单且幂幺，故为恒等。
\end{proof}

\subsubsection{函子 $\mathrm{WD}_{\zeta,\varphi}$}

设 $r\colon W_K\to\GL_n(L)$ 是关于 $L$ 上离散拓扑的连续表示，其中 $L/\mathbb{Q}_p$。
则 $r(\sigma)$ 的特征值（在 $\overline{\mathbb{Q}}_l$ 中）对一切 $\sigma\in W_K$ 都是 $l$-进单位，
当且仅当它们对 $r(\varphi)$（或对任一 $\tau\notin I_K$ 的 $r(\tau)$）是 $l$-进单位，
当且仅当 $\det\rho(\varphi)\in\mathcal{O}_L^{\times}$ 且 $\varphi$ 的特征多项式落在 $\mathcal{O}_L[X]$ 中。
此时称 $r$ 是 \emph{$l$-整的}或 \emph{$l$-有界的}。

\begin{corollary}
固定如上 $\zeta=\{\zeta_n\}$ 与 $\varphi$。则存在范畴等价
\[
\mathrm{WD}=\mathrm{WD}_{\zeta,\varphi}\colon \Rep_L(G_K)\longrightarrow
\mathrm{wd}^{l\text{-}\mathrm{int}}_L(W_K),
\]
从 $L$-向量空间上 $G_K$ 的 $l$-进表示范畴，到 $L$-向量空间上 $W_K$ 的 $l$-整
Weil--Deligne 表示范畴。函子 $\mathrm{WD}$ 由
\[
\mathrm{WD}(\rho)=(\rho',N),\qquad
\rho'(\sigma\varphi^{n}):=\rho(\sigma\varphi^{n})\exp\bigl(-t_l(\sigma)N\bigr)
\]
给出。
\end{corollary}

\begin{remark}
\leavevmode
\begin{enumerate}[label=(\arabic*),leftmargin=2.2em]
\item $\mathrm{WD}_{\zeta',\varphi'}$ 与 $\mathrm{WD}_{\zeta,\varphi}$ 自然同构。
\item $\mathrm{wd}^{l\text{-}\mathrm{int}}_L(W_K)$ 中的 $\varphi$-半单表示，对应于 $\rho(\varphi)$
      为半单算子的 $l$-进表示 $\rho$。
\item 设 $\mathrm{WD}(\rho)=(\rho',N)$，则 $\rho$ 不可约当且仅当 $N=0$ 且 $\rho'$ 不可约。
\end{enumerate}
\end{remark}

由上面两个推论，我们有：

\begin{corollary}
沿用上面的记号，$\rho\colon G_K\to\GL(V)$ 对应 $\mathrm{WD}(\rho)=(\rho',N)$，则下列等价：
\begin{enumerate}[label=(\arabic*),leftmargin=2.2em]
\item $\rho$ 在 $L$ 的某次基变换后半单；
\item $\rho$ 是 $\varphi$-半单的（即 $\varphi$ 半单地作用）且 $\rho(I_K)$ 有限；
\item $\rho'$ 是 $\varphi$-半单的且 $N=0$。
\end{enumerate}
\end{corollary}

\begin{proof}
显然 (1)$\Leftrightarrow$(3)$\Rightarrow$(2)。最后，(2)$\Rightarrow$(3) 由「基本性质」中紧接
Weil--Deligne 表示定义之后的最后一个断言得到；或者也可以这样看出 (2)$\Rightarrow$(1)：
(2) 蕴含在把 $L$ 换成 $\overline{L}$ 后，$\langle\rho(\varphi)\rangle$ 是半单表示，
而 $\langle\rho(\varphi)\rangle<\rho(G_K)$ 是有限指数的，于是可以应用引理 1.1。
\end{proof}

\subsubsection{Weil--Deligne 惯性类型与 Galois 惯性类型}

\begin{definition}
仍设 $V$ 是有限维 $L$-向量空间。
\begin{enumerate}[label=(\arabic*),leftmargin=2.2em]
\item 一个\emph{Weil--Deligne 惯性类型}（有时也简称\emph{惯性类型}）是指表示
\[
\tau\colon I_K\longrightarrow\GL(V),
\]
使得存在 $\varphi$-半单的 Weil--Deligne 表示 $(\rho',N)\colon W_K\to\GL(V)$ 满足
$\tau=\rho'|_{I_K}$。特别地，$\tau$ 的像有限，且不携带关于单值算子 $N$ 的任何信息。
\item 一个\emph{Galois 惯性类型}是指连续的 $\varphi$-半单表示
\[
\xi\colon I_K\longrightarrow\GL(V),
\]
它可以延拓为连续的（$l$-进）表示 $\rho\colon W_K\to\GL(V)$。
\end{enumerate}
\end{definition}

特别地，Galois 惯性类型限制为 Weil--Deligne 惯性类型：取 $\rho$ 为 $\xi$ 的某个延拓，
令 $(\rho',N)=\mathrm{WD}(\rho)$，再取 $\rho'|_{I_K}$。我们下面用 Galois 惯性类型的分类
更直接地描述这件事。

两种惯性类型都可以看作如下数据。$n$ 维（$\varphi$-半单的）Weil--Deligne 表示都形如
$\bigoplus_i\mathrm{Sp}_{m_i}(r_i)$，其中 $r_i$ 是 $W_K$ 的 $n_i$ 维不可约表示，
$\sum_i n_im_i=n$。

\begin{lemma}
设 $(r,N)=\bigoplus_{i=1}^{d}\mathrm{Sp}_{m_i}(r_i)$ 与
$(r',N')=\bigoplus_{i=1}^{d'}\mathrm{Sp}_{m'_i}(r'_i)$ 是 $W_K$ 的两个 $n$ 维
Weil--Deligne 表示，如上（即所有 $r_i$、$r'_i$ 都不可约）。设 $\rho$、$\rho'$ 是
$\mathrm{WD}^{-1}$ 给出的相应 $l$-进表示。
\begin{itemize}[leftmargin=1.6em]
\item 令 $\tau=r|_{I_K}$、$\tau'=r'|_{I_K}$ 是两个 Weil--Deligne 惯性类型；
\item 令 $\xi=\rho|_{I_K}$、$\xi'=\rho'|_{I_K}$ 是两个相应的 Galois 惯性类型。
\end{itemize}
则
\begin{enumerate}[label=(\arabic*),leftmargin=2.2em]
\item 作为 $I_K$ 的表示，$\tau\simeq\tau'$，当且仅当把 $\{r_i\}$（带重复）与 $\{r'_i\}$
      相差非分歧特征标的扭转后两者一致；
\item $\xi\simeq\xi'$，当且仅当 $d=d'$，且适当重排指标后 $m_i=m'_i$、
      $r_i\simeq r'_i\otimes\eta_i$（$\eta_i$ 是 $W_K$ 的非分歧特征标）。
\end{enumerate}
\end{lemma}

\begin{proof}
由引理 4.3 立即得到。
\end{proof}

\begin{center}
\textbf{总结（对应关系）。}
\end{center}

在上述对应与相应的 $\varphi$-半单部分之下，
\[
\Rep_L(G_K)\ \longleftrightarrow\ \mathrm{wd}^{l\text{-}\mathrm{int}}_L(W_K),
\qquad
\Rep^{\varphi\text{-}\mathrm{ss}}_L(G_K)\ \longleftrightarrow\
\mathrm{wd}^{\varphi\text{-}\mathrm{ss},\,l\text{-}\mathrm{int}}_L(W_K),
\]
我们有
\[
\Rep^{\mathrm{ss}}_L(G_K)=\bigl\{\rho\in\Rep^{\varphi\text{-}\mathrm{ss}}_L(G_K):
\rho(I_K)<\infty\bigr\}
\ \longleftrightarrow\
\bigl\{(\rho',N)\in\mathrm{wd}^{\varphi\text{-}\mathrm{ss},\,l\text{-}\mathrm{int}}_L(W_K):N=0\bigr\}.
\]

\begin{keytakeaways}
\begin{itemize}[leftmargin=1.5em]
\item Weil--Deligne 表示 $=$（$W_K$ 的连续表示 $\rho$）$+$（幂零算子 $N$，带 $q^{-\deg}$ 的
      共轭关系）。$N$ 必然幂零，且 $\varphi$-半单情形下 $N=0$ 恰对应「单」。
\item 任何 $\varphi$-半单的 $(\rho,N)$ 都是 $\bigoplus\mathrm{Sp}_m(r)$；$\mathrm{Sp}_m$ 是
      把 $N$ 的链长 $m$ 与 Tate 扭转 $\chi^{-i}$ 捆在一起的那个基本构件。
\item $l$-进表示与 Weil--Deligne 表示由 $\mathrm{WD}_{\zeta,\varphi}$ 等价；
      不可约 $l$-进表示 $\Leftrightarrow$ $N=0$ 且 $\rho'$ 不可约。
\item 惯性类型（Weil--Deligne 型 / Galois 型）只保留限制到 $I_K$ 的信息：
      前者忘掉 $N$，后者保留 Frobenius 的半单作用。第 4 节的惯性 Langlands 对应
      正是在这两类数据上做文章。
\end{itemize}
\end{keytakeaways}
